\documentclass[conference]{IEEEtran}
\usepackage{cite}
\usepackage{amsmath,amssymb,amsfonts}
\usepackage{graphicx}
\usepackage{booktabs}
\usepackage{hyperref}
\usepackage{listings}
\usepackage{xcolor}

\lstset{
    basicstyle=\ttfamily\footnotesize,
    breaklines=true,
    frame=single,
    backgroundcolor=\color{gray!10}
}

\hypersetup{colorlinks=true, linkcolor=blue, citecolor=blue, urlcolor=blue}

\title{SU(3) Matrix Arithmetic over $A_{31}[i]$: \\
       A Deterministic Unitary Coprocessor Extension}

\author{
\IEEEauthorblockN{John Curley}
\IEEEauthorblockA{Independent Researcher, SPU-13 Project\\
\texttt{github.com/pinnacletechnologysolutionsltd/SPU}}
}

\begin{document}

\maketitle

\begin{abstract}
We present an extension of the $A_{31}$ split biquadratic algebra over the
Mersenne prime M31 ($2^{31}-1$) to the degree-8 complex algebra $A_{31}[i]$,
enabling deterministic $3\times3$ unitary matrix arithmetic with zero
floating-point operations. The design reuses the existing M31 multiplier
through a 4-phase time-division multiplexing scheme, running either as a
dedicated standalone sidecar instance or as an integrated sidecar borrowing
a top-level shared instance. On Artix-7, synthesis maps the logical
16-product M31 datapath to 64 DSP48E1 slices. The FSM sequences 27 complex
$A_{31}[i]$ multiplies to complete a full $3\times3$ matrix product in
approximately 224 compute cycles, with row-major result elements emitted as
their accumulators complete. A Python oracle (20 checks) plus multiplier and
sidecar Verilog testbenches verify bit-exact correctness of the arithmetic,
the Gell-Mann structure constants, the conjugate transpose, determinant
checks, and the SPI-visible streaming protocol. The RTL is open-source and
integrated in the Artix-7 SU3 and SU3SHARE spins; this paper reports
simulation, synthesis, Wukong Artix-7 2\,MHz route/bitstream results, JTAG
configuration, and functional SPI/QR silicon verification over the Wukong
J11 header, including all nine dense-matrix result elements on the
shared-multiplier SU3SHARE image.
\end{abstract}

\section{Introduction}

\subsection{Motivation}

Special unitary groups SU($N$) underpin much of modern physics --- lattice
QCD, topological quantum computing, anyonic braid models, and fibre-bundle
gauge theories. In all of these applications, the core computational
primitives are matrix multiplication and coherence checks over the group.
Standard implementations use IEEE-754 floating-point arithmetic, which
accumulates rounding error over long unitary evolution sequences. For
simulations requiring millions of steps --- or for real-time control of
topological quantum systems --- this drift can destroy the structural
invariants that define the group.

\subsection{Approach}

We extend the existing $A_{31}$ split biquadratic algebra over the Mersenne
prime M31 ($2^{31}-1$) by adjoining $\sqrt{-1}$, giving the degree-8 algebra
$A_{31}[i]$. This algebra contains both $\sqrt{3}$ (from $A_{31}$) and
$\sqrt{-1}$ (the complex unit), enabling exact representation of the SU(3)
generators as matrices over $A_{31}[i]$. All arithmetic is exact modular
arithmetic --- no floating point, no rounding, no division in hot paths.

The design instantiates or borrows the same \texttt{spu13\_m31\_multiplier}
module used by the RPLU~v2 pipeline~\cite{rplu2026}, sequencing four
$A_{31}$ base products per complex multiply through a 4-phase time-division
multiplexer. A full $3\times3$ matrix product completes in approximately 224
compute cycles, emitting row-major result elements as each $C[i][j]$
accumulator completes. The standalone module keeps its own M31 multiplier
instance; the SU3SHARE Artix spin instead muxes a top-level shared instance
so the integrated image does not add a second 64-DSP M31 block.

\subsection{What This Paper Is and Is Not}

This paper reports the RTL design, mathematical foundation, and simulation
verification of an SU(3) matrix multiply unit over $A_{31}[i]$. The RTL
compiles under Icarus Verilog and passes all test cases. The Wukong Artix-7
SU3 and SU3SHARE spins synthesize cleanly with the streaming sidecar and
have produced routed 2\,MHz bring-up bitstreams that configure over
DirtyJTAG. The SU3 SPI protocol has been functionally tested on FPGA
hardware: an RP2350 streamed the dense A/B fixture over Wukong J11 and read
exact QR commits for three selected result elements on the standalone
sidecar proof, then for all nine result elements on the SU3SHARE image,
which also preserves the RPLU2 config/QR regression on the same bitstream.
This paper does not claim lattice-QCD-scale capability or timing closure
beyond the 2\,MHz bring-up target.

\section{Mathematical Foundation}

\subsection{$A_{31}$: The Split Biquadratic Algebra over M31}

The base field is the Mersenne prime field $\mathbb{F}_p$ with
$p = 2^{31}-1 = 2147483647$ (M31). Since $p \equiv 3 \pmod 4$, the
polynomial $x^2+1$ is irreducible --- $\sqrt{-1}$ does not exist in the
base field. However, $3^{(p-1)/2} \equiv -1$ and $5^{(p-1)/2} \equiv -1$,
so $\sqrt{3}$ and $\sqrt{5}$ exist in $\mathbb{F}_{p^2}$.

The split biquadratic algebra is
\begin{equation}
A_{31} = \mathbb{F}_p[u,v]\,/\,(u^2-3,\; v^2-5)
\end{equation}
with basis $[1, u, v, uv]$, representing elements as 4-tuples
$(c_0, c_1, c_2, c_3) \in [0, p-1]^4$. The algebra is \emph{split} because
15 is a quadratic residue: $\sqrt{15}$ exists in $\mathbb{F}_p$, so the
extension contains zero divisors. The conjugate reduction
tower~\cite{rplu2026} inverts units in approximately 76 cycles and traps
zero-norm elements via FLAGS.V.

\subsection{Complex Extension: $A_{31}[i]$}

We adjoin $i = \sqrt{-1}$ to $A_{31}$:
\begin{equation}
A_{31}[i] = A_{31}[x]\,/\,(x^2+1).
\end{equation}
An element of $A_{31}[i]$ is a pair $(r, s)$ with $r, s \in A_{31}$,
representing $r + is$. The degree-8 basis over $\mathbb{F}_p$ is
$[1, u, v, uv, i, iu, iv, iuv]$. Multiplication follows the complex rule
\begin{equation}
(r_1 + is_1)(r_2 + is_2) = (r_1 r_2 - s_1 s_2) + i(r_1 s_2 + s_1 r_2),
\end{equation}
so each complex multiply requires four $A_{31}$ base products
($r_1r_2$, $s_1s_2$, $r_1s_2$, $s_1r_2$).

The conjugate reduction tower generalises to three stages
($\mathbb{F}_p \to \mathbb{F}_p(\sqrt{3}) \to \mathbb{F}_p(\sqrt{3},\sqrt{5})
\to \mathbb{F}_p(\sqrt{3},\sqrt{5},\sqrt{-1})$), giving an inversion latency
of approximately 114 cycles for $A_{31}[i]$.

\subsection{SU(3) Structure Constants in $A_{31}[i]$}

The eight Gell-Mann matrices $\lambda_1,\dots,\lambda_8$ generate
$\mathfrak{su}(3)$~\cite{gellmann1962}. Their entries lie in
$\{0, \pm1, \pm\tfrac{1}{2}, \pm\tfrac{\sqrt{3}}{2}\}$. In $A_{31}[i]$:
\begin{itemize}
\item $\pm1$ and $\pm\tfrac{1}{2}$ are field elements
      ($2^{-1} = 1073741824$ in M31);
\item $\sqrt{3}$ exists natively as $u$;
\item $\pm\tfrac{\sqrt{3}}{2} = \pm\tfrac{1}{2}\cdot\sqrt{3}$ is a single
      $A_{31}$ multiplication;
\item $i$ is the complex unit, present as the imaginary basis element.
\end{itemize}
For this work we verify that $\lambda_1^2 = \operatorname{diag}(1,1,0)$,
$\lambda_3$ is diagonal with entries $(1,-1,0)$, $\lambda_8$ carries the
$\sqrt{3}$ dependence correctly, and matrix multiplication over $A_{31}[i]$
is closed for $3\times3$ matrices. The determinant of every SU(3) element
equals $1 + 0i$ in $A_{31}[i]$.

\section{RTL Design}

\subsection{Multiplier Topography}

The \texttt{spu13\_su3\_mult} module uses the 2-stage pipelined
\texttt{spu13\_m31\_multiplier} --- the same multiplier design used by the
RPLU~v2 pipeline. Each $A_{31}$ product (4 components, 16 logical parallel
$32\times32$ integer products) completes with 2-cycle throughput; Artix-7
synthesis maps the logical products to 64 DSP48E1 slices. A complex
$A_{31}[i]$ multiply is sequenced as four phases through this multiplier
(Table~\ref{tab:phases}).

\begin{table}[ht]
\centering
\caption{4-phase TDM sequence per complex multiply.}
\label{tab:phases}
\begin{tabular}{clc}
\toprule
Phase & Operation & Cycles \\
\midrule
0 & $RR = \mathrm{real}(A)\cdot\mathrm{real}(B)$ & 2 \\
1 & $II = \mathrm{imag}(A)\cdot\mathrm{imag}(B)$ & 2 \\
2 & $RI = \mathrm{real}(A)\cdot\mathrm{imag}(B)$ & 2 \\
3 & $IR = \mathrm{imag}(A)\cdot\mathrm{real}(B)$ & 2 \\
\bottomrule
\end{tabular}
\end{table}

After all four phases the accumulators update as
$c_{\mathrm{real}} \mathrel{+}= RR - II$ and
$c_{\mathrm{imag}} \mathrel{+}= RI + IR$ ($A_{31}$ subtraction and
addition). Each complex multiply takes 8 cycles total, fully pipelined: the
multiplier starts a new $A_{31}$ product every 2 cycles while the previous
result is being accumulated.

\subsection{$3\times3$ Matrix Multiply FSM}

The FSM implements
\begin{equation}
C[i][j] = \sum_{k=0}^{2} A[i][k]\cdot B[k][j], \qquad i,j \in \{0,1,2\}.
\end{equation}
The interface uses element-wise loading: the host writes 9 elements
(256 bits each) into matrix A via \texttt{load\_a} strobes, then 9 into
matrix B via \texttt{load\_b} strobes. Each element packs the four
imaginary then four real 32-bit components. The FSM sequences 27 complex
multiplies in row-major order; each result element is emitted as soon as
its three $k$-terms are accumulated, so the result stream is part of the
compute schedule rather than a separate drain stage. Total latency is
approximately 224 compute cycles.

\subsection{Resource Utilisation}

\begin{table}[ht]
\centering
\caption{Resource and timing snapshot (Wukong Artix-7 100T).}
\label{tab:resources}
\begin{tabular}{lc}
\toprule
Item & Value \\
\midrule
DSP slices & 64 DSP48E1 \\
Estimated LCs & $\sim$6.2k multiplier, $\sim$8.2k SU3 spin \\
Packed FFs (SU3 spin) & 9,488 \\
Packed LUT cells (SU3 spin) & 21,092 (16\% of XC7A100T) \\
SU3SHARE integration snapshot & 60,837 LUTX, 16,478 FFX, 64 DSP \\
Post-route timing (SU3) & 51.58\,MHz max, PASS at 2\,MHz target \\
Post-route timing (SU3SHARE) & 3.67\,MHz max, PASS at 2\,MHz target \\
Cycles per $3\times3$ multiply & $\sim$224 \\
Projected time at 200\,MHz & $\sim$1.2\,$\mu$s \\
\bottomrule
\end{tabular}
\end{table}

The sidecar is parametrised on \texttt{EXTERNAL\_MULT}: the default
instantiates a dedicated multiplier, while the SU3SHARE spin sets it to
borrow one top-level shared \texttt{spu13\_m31\_multiplier} so the
integrated image does not duplicate the 64-DSP block.

\subsection{Worked Example: $\lambda_1 \times \lambda_1$}

Consider the Gell-Mann generator $\lambda_1$ multiplied by itself:
\begin{equation}
\lambda_1 = \begin{pmatrix} 0&1&0\\ 1&0&0\\ 0&0&0 \end{pmatrix},
\qquad
\lambda_1^2 = \begin{pmatrix} 1&0&0\\ 0&1&0\\ 0&0&0 \end{pmatrix}.
\end{equation}
Each matrix element is an $A_{31}[i]$ value: a pair of 4-tuples
(real, imaginary) over M31. Here all values are real with only $c_0$
components equal to 0 or 1. The computation
$C[0][0] = A[0][0]B[0][0] + A[0][1]B[1][0] + A[0][2]B[2][0]$ evaluates to
$0\cdot0 + 1\cdot1 + 0\cdot0 = 1$. The middle term requires a full
$A_{31}[i]$ multiply: the 4-phase TDM sequence produces $RR=1$, $II=0$,
$RI=0$, $IR=0$, giving $c_{\mathrm{real}}=1$, $c_{\mathrm{imag}}=0$, so
$C[0][0] = (1,0,0,0)$. This example is verified by both the Python oracle
and the Verilog testbench.

\section{Verification}

\subsection{Python Oracle}

The oracle (\texttt{test\_su3\_oracle.py}) implements $A_{31}[i]$
arithmetic, the Gell-Mann matrices, $3\times3$ matrix multiplication, the
conjugate transpose, determinant, and coherence checks. Twenty checks
cover: $A_{31}$ arithmetic (3), $A_{31}[i]$ arithmetic including
$i^2 = -1$ (2), matrix-multiply identity and closure (2), the dense packed
fixture shared with the RTL testbench (1), Gell-Mann structure (8), product
closure (1), self-adjoint coherence
$\lambda_1\lambda_1^\dagger = \lambda_1^2$ (1), and determinant of the
identity (2). All 20 checks pass. The oracle uses standard Python integers
reduced modulo M31 --- no floating point, no external libraries.

\subsection{RTL Testbench}

The Verilog testbench (\texttt{spu13\_su3\_mult\_tb.v}) loads test
matrices through the element-wise interface and asserts bit-exact output
against the public result stream across five cases: identity closure,
$\lambda_1 \times \lambda_1 = \operatorname{diag}(1,1,0)$, right-operand
preservation, zero propagation with accumulator clearing, and a dense
$A_{31}[i]$ product exercising all lanes and three-term accumulation with
oracle-derived packed constants. All cases pass under Icarus Verilog, and
the bench is auto-discovered by the project regression suite. The Python
oracle serves as the golden reference, mirroring the verification
methodology of the Lucas MAC~\cite{lucas2026} and RPLU~v2~\cite{rplu2026}
pipelines.

\subsection{Wukong Artix-7 Silicon Verification}

On 2026-07-04, the sidecar-only SU3 bitstream was SRAM-loaded through
RP2040 DirtyJTAG at 1\,MHz. RP2350 smoke firmware streamed the dense A/B
matrix fixture over the Wukong J11 SPI path (command \texttt{0xB1}) at
100\,kHz with 20\,$\mu$s CS-setup, read-turnaround, CRC-hold, and
CS-recovery guard delays. Three independent result elements, read back
through the QR commit path, matched the oracle exactly
(Table~\ref{tab:silicon}); the capture ended with \texttt{SU3\_J11: PASS},
and a 40-second capture completed thirteen three-case passes. A 5\,$\mu$s
guard-delay probe produced an intermittent invalid QR read, so 20\,$\mu$s
is the current practical margin.

On 2026-07-06, the SU3SHARE smoke was expanded to all nine dense-product
result elements, read through QR lanes 0--8. Two complete capture loops
reported exact matches for every element, and the same bitstream then
passed the RPLU2 config/QR regression.

\begin{table}[ht]
\centering
\caption{Silicon QR readback, standalone SU3 spin (2026-07-04).}
\label{tab:silicon}
\resizebox{\columnwidth}{!}{%
\begin{tabular}{ccll}
\toprule
Result & QR lane & A & B \\
\midrule
0 & 2 & \texttt{7FFE271F7FFC43EF} & \texttt{7FFF6B677FFED36F} \\
4 & 5 & \texttt{7FFD2A6B7FFA47FF} & \texttt{7FFF196B7FFE2E9F} \\
8 & 8 & \texttt{7FFBF6DF7FF7DE5F} & \texttt{7FFEB5277FFD653F} \\
\midrule
Result & QR lane & C & D \\
\midrule
0 & 2 & \texttt{00021510000446A0} & \texttt{0000A30000014F30} \\
4 & 5 & \texttt{00034B480006BAA8} & \texttt{00010678000218A8} \\
8 & 8 & \texttt{0004CAA00009C0F0} & \texttt{00018250000312E0} \\
\bottomrule
\end{tabular}
}
\end{table}

\section{Integration Path}

The SU(3) module is part of the Artix-7 synthesis script; Yosys prunes it
from builds that do not instantiate it, so it adds zero cost until enabled.
The element-wise load interface rides the existing SPI command path
(\texttt{0xB1}): \texttt{EA} starts and selects a result element,
\texttt{E8}/\texttt{E9} stream A/B element chunks, and \texttt{EB} commits
the captured 256-bit result through QR A/B/C/D. A full $3\times3$ multiply
requires 144 chunk writes (9 elements each for A and B), one start command,
and one read command per selected result element.

\section{Discussion}

\subsection{Applications}

The primary application is not large-scale lattice QCD --- the 32-bit M31
field is too small for double-precision physics. The natural niche is
topological and Fibonacci anyon models~\cite{freedman2000} where the
algebra is naturally small-domain, exactness matters more than precision,
and the SU(3) structure maps onto fusion-rule bookkeeping. A secondary
application is exact unitary benchmarking: verifying that a sequence of
unitary operations closes to the identity within a finite algebra, without
floating-point drift.

\subsection{Limitations}

\begin{itemize}
\item \textbf{Guarded bring-up link:} the silicon proof uses a 100\,kHz SPI
stream with 20\,$\mu$s guard delays; a 5\,$\mu$s probe showed intermittent
QR read invalidity, so link tuning should proceed from the 20\,$\mu$s
baseline.
\item \textbf{Small field:} M31 is a 32-bit prime; operations are exact but
low-precision. Larger applications require multiple residues or larger
Mersenne primes.
\item \textbf{Multiplication only:} group exponentiation, logarithm, and
generator decomposition are not implemented.
\item \textbf{No exponential map:} $\exp(X)$ for
$X \in \mathfrak{su}(3)$ involves factorials that are not well-defined in
finite characteristic; the design works entirely in the group domain via
direct matrix multiplication.
\end{itemize}

\subsection{Future Work}

Generator exponentiation via repeated squaring; SU($N$) generalisation for
$N > 3$; integration with the Lucas MAC for combined $\phi$ + unitary
arithmetic; a standalone SU(3) probe on Tang 25K following the PHSLK
microprobe pattern; and a multi-residue (CRT) extension for
higher-precision unitary arithmetic.

\section*{Appendix: Python Oracle (Abridged)}

\begin{lstlisting}[language=Python]
P = 2147483647  # M31

def a31_mul(a, b):
    """A31 multiplication: 16 products mod M31."""
    c0, c1, c2, c3 = a; d0, d1, d2, d3 = b
    return (
        (c0*d0 + 3*c1*d1 + 5*c2*d2 + 15*c3*d3) % P,
        (c0*d1 + c1*d0 + 5*c2*d3 + 5*c3*d2) % P,
        (c0*d2 + c2*d0 + 3*c1*d3 + 3*c3*d1) % P,
        (c0*d3 + c1*d2 + c2*d1 + c3*d0) % P,
    )

def ca_mul(r1, s1, r2, s2):
    """A31[i] multiply: (r1 + i s1)(r2 + i s2)."""
    rr = a31_mul(r1, r2); ss = a31_mul(s1, s2)
    rs = a31_mul(r1, s2); sr = a31_mul(s1, r2)
    return (a31_sub(rr, ss), a31_add(rs, sr))
\end{lstlisting}

Full oracle with 20 checks and Gell-Mann matrix definitions:
\texttt{software/tests/test\_su3\_oracle.py}.

\begin{thebibliography}{9}

\bibitem{gellmann1962}
M.~Gell-Mann, ``Symmetries of baryons and mesons,''
\emph{Phys. Rev.}, vol.~125, p.~1067, 1962.

\bibitem{freedman2000}
M.~Freedman, M.~Larsen, and Z.~Wang, ``A modular functor which is universal
for quantum computation,'' arXiv:quant-ph/0001108, 2000.

\bibitem{crandall2005}
R.~E.~Crandall and C.~Pomerance, \emph{Prime Numbers: A Computational
Perspective}, 2nd~ed. Springer, 2005.

\bibitem{lucas2026}
J.~Curley, ``Zero-drift Lucas-prime phinary arithmetic for exact quantum
phase coherence on FPGA,'' 2026, \emph{(in preparation)}.

\bibitem{rplu2026}
J.~Curley, ``A hardware jet algebra coprocessor over a split M31
biquadratic algebra,'' 2026, \emph{(in preparation)}.

\bibitem{spu13repo}
SPU-13 Project, open-source RTL and oracles.
\url{https://github.com/pinnacletechnologysolutionsltd/SPU}

\end{thebibliography}

\end{document}
